\documentclass[11pt,a4paper]{article}

\usepackage[utf8]{inputenc}
\usepackage[T1]{fontenc}
\usepackage[spanish,es-nodecimaldot]{babel}
\usepackage{lmodern}
\usepackage{microtype}
\usepackage{geometry}
\usepackage{enumitem}
\usepackage{booktabs}
\usepackage{tabularx}
\usepackage{xcolor}
\usepackage{tikz}
\usepackage[hidelinks]{hyperref}
\usetikzlibrary{arrows.meta,positioning}

\geometry{top=2.2cm,bottom=2.2cm,left=2.5cm,right=2.5cm}
\setlength{\parindent}{0pt}
\setlength{\parskip}{0.55em}
\setlist[enumerate]{leftmargin=*,itemsep=0.35em,topsep=0.35em}
\renewcommand{\arraystretch}{1.2}

\newcommand{\autor}{Francisco Clarke}
\newcommand{\lu}{120547}
\newcommand{\director}{Ana Gabriela Maguitman}
\newcommand{\codirector}{Federico Mart\'in Schmidt}

\begin{document}

\begin{center}
  {\large\textbf{UNIVERSIDAD NACIONAL DEL SUR}}\\
  {\large Departamento de Ciencias e Ingenier\'ia de la Computaci\'on}\\[0.8em]
  {\large\textbf{TRABAJO FINAL DE INGENIER\'IA EN SISTEMAS DE INFORMACI\'ON}}\\
  {\Large\textbf{PLAN DE TRABAJO}}\\[1em]

  {\large\textbf{Dise\~no e implementaci\'on de una plataforma de
  orquestaci\'on de agentes de lenguaje para el desarrollo de software y
  an\'alisis exploratorio de una pol\'itica de descomposici\'on adaptativa}}\\[1em]

  \begin{tabular}{rl}
    Alumno: & \autor \\
    LU: & \lu \\
    Directora: & \director \\
    Codirector: & \codirector
  \end{tabular}
\end{center}

\section*{1. Objetivos}

Los agentes basados en modelos de lenguaje han demostrado capacidad para
interactuar con repositorios de software, modificar archivos y ejecutar
herramientas de desarrollo. Trabajos como SWE-agent muestran que la interfaz
provista al agente y la forma en que se estructura su interacci\'on con el
entorno influyen en su desempe\~no \cite{yang2024sweagent}. Otros enfoques, como
ChatDev y MetaGPT, exploran la colaboraci\'on entre agentes especializados para
cubrir distintas etapas del desarrollo de software
\cite{qian2024chatdev,hong2024metagpt}.

Sin embargo, una funcionalidad de software de alcance medio o grande no se
reduce a una \'unica tarea de generaci\'on de c\'odigo. Es necesario interpretar
el objetivo, dividirlo en unidades ejecutables, explicitar dependencias y
contratos, coordinar trabajo potencialmente paralelo, verificar los cambios y
recomponer los resultados sin perder trazabilidad. La ejecuci\'on paralela
basada en grafos de dependencias ya ha mostrado ventajas de latencia y costo en
otros dominios de orquestaci\'on de herramientas con modelos de lenguaje
\cite{kim2024llmcompiler}; su aplicaci\'on al desarrollo de software agrega
desaf\'ios vinculados con el estado compartido del repositorio, la validaci\'on
del c\'odigo y la integraci\'on sem\'antica de cambios.

Este trabajo propone el dise\~no, implementaci\'on y an\'alisis de una
plataforma denominada \textbf{Conductor}, que transforma una especificaci\'on en
lenguaje natural en un proceso de desarrollo supervisable. Su arquitectura se
organiza alrededor de tres mecanismos principales: descomposici\'on recursiva,
ejecuci\'on aislada e integraci\'on incremental. Cada mecanismo constituye un
problema de dise\~no propio y su articulaci\'on permite recorrer el ciclo completo
desde la intenci\'on inicial hasta la obtenci\'on de cambios integrados y
verificados.

La Figura~\ref{fig:arquitectura-conductor} resume este flujo y los mecanismos
transversales de supervisi\'on y persistencia que permiten inspeccionarlo y
reanudarlo.

\begin{figure}[ht]
  \centering
  \begin{tikzpicture}[
    node distance=5mm,
    flow/.style={
      draw=black!55,
      rounded corners=2pt,
      fill=black!4,
      align=center,
      text width=2.65cm,
      minimum height=1.45cm,
      font=\small
    },
    endpoint/.style={
      draw=black!55,
      rounded corners=2pt,
      fill=blue!5,
      align=center,
      text width=2.25cm,
      minimum height=1.1cm,
      font=\small
    },
    control/.style={
      draw=black!45,
      dashed,
      rounded corners=2pt,
      fill=orange!7,
      align=center,
      text width=8.9cm,
      minimum height=1cm,
      font=\small
    },
    arrow/.style={-{Stealth[length=2mm]},thick,draw=black!65},
    controlarrow/.style={-{Stealth[length=1.7mm]},dashed,draw=black!50}
  ]
    \node[endpoint] (input) {Especificaci\'on\\en lenguaje natural};
    \node[flow,right=of input] (planning)
      {\textbf{Descomposici\'on}\\DAG, contratos e interfaces};
    \node[flow,right=of planning] (execution)
      {\textbf{Ejecuci\'on}\\worktrees, agentes y validaci\'on};
    \node[flow,right=of execution] (integration)
      {\textbf{Integraci\'on}\\composici\'on y reparaci\'on};
    \node[endpoint,right=of integration] (result)
      {Resultado\\integrado y verificable};

    \draw[arrow] (input) -- (planning);
    \draw[arrow] (planning) -- (execution);
    \draw[arrow] (execution) -- (integration);
    \draw[arrow] (integration) -- (result);

    \node[control,below=8mm of execution] (supervision)
      {\textbf{Supervisi\'on y estado durable:} revisi\'on humana, eventos,
      decisiones, checkpoints y recuperaci\'on};
    \draw[controlarrow] (supervision.north west) -- (planning.south);
    \draw[controlarrow] (supervision.north) -- (execution.south);
    \draw[controlarrow] (supervision.north east) -- (integration.south);
  \end{tikzpicture}
  \caption{Arquitectura conceptual de Conductor.}
  \label{fig:arquitectura-conductor}
\end{figure}

\subsection*{Objetivo general}

Dise\~nar, implementar y analizar Conductor, una plataforma de orquestaci\'on de
agentes basados en modelos de lenguaje para el desarrollo de software, capaz de
transformar una especificaci\'on en lenguaje natural en un conjunto estructurado
de tareas, ejecutar dichas tareas en entornos aislados e integrar sus resultados
de manera verificable, trazable y supervisable. Como parte del an\'alisis, se
estudiar\'a de forma exploratoria una pol\'itica adaptativa para determinar la
granularidad de la descomposici\'on seg\'un las caracter\'isticas de cada tarea.

\subsection*{Objetivos espec\'ificos}

\begin{enumerate}
  \item Definir una arquitectura de planificaci\'on basada en un grafo jer\'arquico
  e implementar un mecanismo de descomposici\'on recursiva que produzca tareas,
  dependencias, criterios de aceptaci\'on e interfaces compartidas.

  \item Implementar una ejecuci\'on aislada que asigne tareas a agentes, controle
  el alcance de sus modificaciones, ejecute validaciones y registre evidencia
  de los cambios realizados.

  \item Implementar una integraci\'on incremental que recomponga los resultados
  seg\'un el grafo, conserve su trazabilidad y permita reparar conflictos a partir
  del objetivo y los contratos definidos durante la descomposici\'on.

  \item Proporcionar mecanismos de supervisi\'on y recuperaci\'on para observar el
  progreso, registrar decisiones, intervenir ante fallas y reanudar ejecuciones.

  \item Dise\~nar una pol\'itica adaptativa de granularidad que considere la
  complejidad de cada tarea, informaci\'on estructural del repositorio y el costo
  potencial de nuevas divisiones.

  \item Analizar el funcionamiento de la plataforma mediante casos seleccionados
  de desarrollo de software y comparar descriptivamente los planes producidos
  por distintas configuraciones de granularidad.
\end{enumerate}

\section*{2. Metodolog\'ia}

El trabajo combinar\'a dise\~no de arquitectura, desarrollo incremental,
verificaci\'on automatizada y an\'alisis exploratorio. La implementaci\'on se
organizar\'a en m\'odulos con responsabilidades expl\'icitas y contratos entre
ellos. El estado persistido y los eventos producidos por cada ejecuci\'on se
utilizar\'an tanto para supervisar el producto como para reconstruir las
decisiones adoptadas en los casos estudiados.

\subsection*{2.1. Descomposici\'on}

La entrada del sistema ser\'a una descripci\'on de una funcionalidad sobre un
repositorio de software. Un agente planificador analizar\'a recursivamente cada
nodo y decidir\'a si representa una unidad cohesiva y verificable o si requiere
una nueva divisi\'on. El resultado ser\'a un grafo dirigido ac\'iclico (DAG) con
dependencias, criterios de aceptaci\'on, alcance esperado e interfaces
compartidas.

Las interfaces actuar\'an como contratos entre tareas: permitir\'an que agentes
aislados trabajen sobre una frontera com\'un sin necesitar los resultados
concretos de sus pares. Las salidas del planificador se validar\'an de forma
estructural y sem\'antica antes de habilitar la ejecuci\'on.

\subsection*{2.2. Ejecuci\'on}

Cada tarea hoja se ejecutar\'a en un \emph{worktree} de Git independiente. El
agente recibir\'a su objetivo, los contratos que consume o produce, el contexto
relevante del repositorio y las validaciones esperadas. Una vez finalizada la
tarea, el sistema utilizar\'a el estado real del repositorio para determinar los
cambios realizados, comprobar que respeten el alcance autorizado y ejecutar
pruebas o comandos de validaci\'on.

Las tareas sin dependencias incompatibles podr\'an ejecutarse en paralelo. La
selecci\'on de cada grupo de ejecuci\'on considerar\'a dependencias, superposici\'on
de archivos y se\~nales de riesgo estructural. Ante una validaci\'on fallida, el
agente podr\'a recibir el diagn\'ostico y realizar un intento acotado de
reparaci\'on.

\subsection*{2.3. Integraci\'on y supervisi\'on}

Los resultados se integrar\'an de abajo hacia arriba siguiendo la jerarqu\'ia del
grafo. Cada contribuci\'on conservar\'a un commit identificable y ser\'a aplicada
al nodo compuesto correspondiente. Cuando aparezca un conflicto, un agente de
integraci\'on recibir\'a el objetivo del nodo padre, las interfaces compartidas,
la intenci\'on de los hijos y los cambios involucrados. La reparaci\'on deber\'a
respetar esos contratos y superar las validaciones del conjunto.

La plataforma persistir\'a el estado del proceso, sus eventos y las decisiones
humanas. Una interfaz web permitir\'a revisar el plan, observar la evoluci\'on del
grafo, responder preguntas y resolver situaciones que no puedan tratarse de
manera aut\'onoma.

\subsection*{2.4. Pol\'itica adaptativa y an\'alisis exploratorio}

La granularidad, el modelo y el esfuerzo de razonamiento son decisiones de
asignaci\'on de recursos. Las interfaces actuales permiten
configurar los dos \~ultimos \cite{openai2025reasoning,anthropic2026adaptive}, y
distintos trabajos estudian pol\'iticas para equilibrar costo y calidad
\cite{chen2023frugalgpt,ong2025routellm}. En Conductor, la granularidad es una
decisi\'on previa: define las tareas y condiciona su ejecuci\'on, paralelismo e
integraci\'on.

Se implementar\'a una pol\'itica que combine una estimaci\'on determinista de
complejidad con una r\'ubrica sensible al costo de dividir. Se considerar\'an el
alcance, las interfaces, los archivos esperados y se\~nales del repositorio,
como s\'imbolos e imports, frente al costo adicional de coordinaci\'on,
integraci\'on y posibles conflictos. Cada recomendaci\'on conservar\'a evidencia
auditable.

El an\'alisis exploratorio se organizar\'a en dos l\'ineas complementarias. La
primera estudiar\'a tres casos completos: una funcionalidad simple, otra de
complejidad intermedia y una con solapamientos
o conflictos esperables. Si los recursos lo permiten, se incorporar\'a un cuarto
caso centrado en la recuperaci\'on ante una falla. Para cada caso se describir\'an
el DAG producido, las interfaces y dependencias, las waves de ejecuci\'on, las
validaciones, los conflictos, las reparaciones, las intervenciones humanas y el
resultado integrado. Tambi\'en se registrar\'an la duraci\'on y el consumo cuando
puedan obtenerse de forma confiable.

La segunda l\'inea analizar\'a la granularidad sin ejecutar las tareas. Sobre un
conjunto acotado de funcionalidades se comparar\'an los planes generados por la
pol\'itica propuesta, la r\'ubrica adaptativa inicial y una granularidad fija
intermedia. Se observar\'an la cantidad total de nodos, la profundidad y el ancho
del DAG, las tareas at\'omicas y compuestas, las interfaces, las dependencias y
los posibles solapamientos entre alcances.

Dentro de cada comparaci\'on se mantendr\'an constantes el repositorio, la
funcionalidad y el modelo. Los resultados se interpretar\'an de manera
descriptiva y cualitativa, sin pretensi\'on de generalizaci\'on estad\'istica. La
selecci\'on autom\'atica de modelos y esfuerzo quedar\'a como trabajo futuro, y la
comparaci\'on con un agente \'unico quedar\'a fuera del alcance.

\section*{3. Cronograma de actividades}

Se adopta un cronograma tentativo de aproximadamente cuatro meses, iniciado
durante la primera mitad de abril de 2026.

\begin{center}
\small
\begin{tabularx}{\textwidth}{@{}Xccccc@{}}
  \toprule
  \textbf{Actividad} & \textbf{Abr.} & \textbf{May.} & \textbf{Jun.} & \textbf{Jul.} & \textbf{Ago.} \\
  \midrule
  1. Relevamiento y definici\'on de arquitectura
    & \(\bullet\) & \(\bullet\) & & & \\
  2. Descomposici\'on y contratos
    & \(\bullet\) & \(\bullet\) & \(\bullet\) & & \\
  3. Ejecuci\'on aislada y validaci\'on
    & & \(\bullet\) & \(\bullet\) & \(\bullet\) & \\
  4. Integraci\'on, persistencia y supervisi\'on
    & & & \(\bullet\) & \(\bullet\) & \\
  5. Pol\'itica adaptativa de granularidad
    & & & & \(\bullet\) & \(\bullet\) \\
  6. Estudio de casos y an\'alisis exploratorio
    & & & & \(\bullet\) & \(\bullet\) \\
  7. Documentaci\'on y redacci\'on de la tesina
    & & \(\bullet\) & \(\bullet\) & \(\bullet\) & \(\bullet\) \\
  \bottomrule
\end{tabularx}
\end{center}

\section*{4. Resultados esperados}

Se espera obtener una plataforma funcional capaz de transformar un objetivo de
desarrollo en un plan estructurado, ejecutar sus tareas en entornos aislados e
integrar los cambios producidos. El sistema deber\'a conservar evidencia de cada
etapa y ofrecer puntos de supervisi\'on y recuperaci\'on ante decisiones ambiguas
o fallas.

El principal resultado t\'ecnico ser\'a la articulaci\'on de la descomposici\'on,
la ejecuci\'on y la integraci\'on en un producto completo. La tesina documentar\'a
la arquitectura, los contratos entre estos mecanismos, las decisiones de
dise\~no adoptadas y las limitaciones identificadas durante su implementaci\'on.

El estudio de casos permitir\'a describir el comportamiento del sistema completo
en escenarios de distinta complejidad, incluyendo su respuesta ante conflictos,
reparaciones o decisiones humanas. El an\'alisis de planificaci\'on documentar\'a
c\'omo var\'ia la estructura del DAG entre configuraciones y aportar\'a evidencia
sobre la utilidad y los l\'imites de la pol\'itica adaptativa, sin formular
generalizaciones estad\'isticas.

\section*{5. Vinculaci\'on con actividades de investigaci\'on o transferencia}

El presente proyecto final no se encuentra vinculado formalmente a ninguna
actividad de investigaci\'on, vinculaci\'on o extensi\'on vigente en el
Departamento de Ciencias e Ingenier\'ia de la Computaci\'on (DCIC).

\section*{6. Firmas}

\vspace{1.8cm}
\begin{center}
  \begin{tabularx}{\textwidth}{@{}>{\centering\arraybackslash}X >{\centering\arraybackslash}X >{\centering\arraybackslash}X@{}}
    \rule{4.2cm}{0.4pt} & \rule{4.2cm}{0.4pt} & \rule{4.2cm}{0.4pt} \\
    \autor & \director & \codirector \\
    Alumno & Directora & Codirector
  \end{tabularx}
\end{center}

\clearpage
\bibliographystyle{ieeetr}
\bibliography{referencias}

\end{document}
